\documentclass[sigconf]{acmart}

\acmConference[CONF '26]{ACM Conference}{June 01--03, 2026}{New York, NY}
\acmDOI{10.1145/1122334.4455667}
\keywords{LaTeX, parsing, editing, agents}

\title{A Surgical LaTeX Editor Agent}

\begin{document}

\begin{abstract}
We present a backend agent that parses LaTeX papers into structural zones
and edits individual sections under natural-language instructions.
\end{abstract}

\maketitle

\section{Introduction}
Academic writing in LaTeX is a slow, error-prone craft. Authors spend
significant effort rephrasing sections, tightening prose, and reconciling
citations across drafts. Existing editors help with syntax but offer little
semantic assistance \cite{Knuth1984Literate, Lamport1994LaTeX}.
In this paper we present a backend agent that parses LaTeX papers into
structural zones and lets the user rewrite individual sections through
natural-language instructions, with surgical edits that never touch the
surrounding document.

\section{Background}
\label{sec:background}
Prior work on academic-writing assistants is reviewed here.

\subsection{Editors}
Existing LaTeX editors focus on syntax, not semantics.

\section{Methodology}
% TODO: write this section

\section{Results}

\section{Conclusion}
We have described a surgical editor for ACM-style papers.

\begin{acks}
We thank the IntelliDraft team for valuable feedback.
\end{acks}

\bibliographystyle{ACM-Reference-Format}
\bibliography{refs}

\end{document}
